\leanverified{paper\_leyang\_haar\_pushforward\_density}
\begin{theorem}[Haar 推前到 Lee--Yang 奇异环的显式密度]\label{thm:leyang-haar-pushforward-density}
令 \(m_{\TT}\) 为单位圆 \(\TT\) 上的归一化 Haar 测度，并令随机变量 \(U\sim m_{\TT}\)。
定义 \(T:=J(U)\)，其中 \(J\) 由 \eqref{eq:lee-yang-J-def} 给出。
则 \(T\) 几乎处处取值于 Lee--Yang 奇异环 \((-\infty,-\tfrac14]\)，且其分布对 Lebesgue 测度绝对连续，密度为
\begin{equation}\label{eq:leyang-haar-pushforward-density}
\boxed{\;
f_T(t)
=\frac{1}{\pi\,|t|\,\sqrt{|1+4t|}}\;\ind_{(-\infty,-1/4]}(t)\; }.
\end{equation}
并且 \(\int_{\RR} f_T(t)\,\dd t=1\)。
\end{theorem}

\begin{proof}
取 \(z\in\TT\) 使 \(U=z^2\)。
由命题 \ref{prop:jouwkowsky-inverse-square}，
\[
T=J(z^2)=-\frac{1}{(z+z^{-1})^2}.
\]
在 \(m_{\TT}\) 下可取 \(z=e^{\mathrm{i}\phi}\) 且 \(\phi\sim\mathrm{Unif}[0,\pi)\)，从而
\[
T=-\frac{1}{4\cos^2\phi}.
\]
置 \(s:=-t\ge \tfrac14\)，并对变换 \(s=\frac{1}{4\cos^2\phi}\) 进行换元计算可得
\(
f_T(t)=\frac{1}{\pi(-t)\sqrt{-4t-1}}
\)
于 \(t\le -\tfrac14\)，等价于 \eqref{eq:leyang-haar-pushforward-density}。
归一化由同一换元直接核验。证毕。
\end{proof}

\begin{corollary}[余弦极坐标族的推前分布刚性]\label{cor:leyang-polar-family-pushforward-rigidity}
\leanverified{paper\_leyang\_polar\_family\_pushforward\_rigidity}
设 \(\gcd(n,d)=1\)，并沿用 \eqref{eq:leyang-polar-family-t-def} 的记号：
\(\alpha=e^{\mathrm{i}\theta/d}\in\TT\)，以及 \(t(\theta):=J(\alpha^{2n})\)。
若 \(\Theta\) 在 \([0,2\pi d)\) 上均匀分布，则随机变量 \(t(\Theta)\) 的分布与 \((n,d)\) 无关，且与定理 \ref{thm:leyang-haar-pushforward-density} 的 \(T\) 同分布；因此其密度亦由 \eqref{eq:leyang-haar-pushforward-density} 给出。
\end{corollary}

\begin{proof}
令 \(\varphi:=\frac{n}{d}\Theta\)。
由于 \(\gcd(n,d)=1\)，映射 \(\Theta\mapsto\varphi\) 将 \([0,2\pi d)\) 的均匀测度推为 \([0,2\pi n)\) 的均匀测度。
由 \eqref{eq:leyang-polar-family-t-closed}，
\(
t(\Theta)=-\frac{1}{4\cos^2\varphi}
\)。
又因 \(\cos^2\) 以周期 \(\pi\) 折叠，故 \(t(\Theta)\) 与 \(-\frac{1}{4\cos^2\Phi}\)（\(\Phi\sim\mathrm{Unif}[0,\pi)\)）同分布。
后者与 \(U=e^{2\mathrm{i}\Phi}\sim m_{\TT}\) 下的 \(J(U)\) 同分布，即为定理 \ref{thm:leyang-haar-pushforward-density} 的结论。证毕。
\end{proof}

\begin{remark}\label{rem:leyang-polar-family-statistical-invariant}
推前坐标 \(t=-\frac{1}{4r(\theta)^2}\) 在概率口径下不依赖于 \((n,d)\)：
参数 \(\frac{n}{d}\) 仅改变轨迹在复平面中的花瓣数与自交结构，而 \(t(\Theta)\) 的推前分布保持不变。
\end{remark}

\endinput
